\documentclass[aps,prb,onecolumn,nofootinbib,superscriptaddress]{revtex4-2}

\usepackage{amsmath,amssymb,amsthm,mathtools,bm}
\usepackage{booktabs}
\usepackage[colorlinks=true,linkcolor=blue,citecolor=blue,urlcolor=blue]{hyperref}
\usepackage{microtype}

\newcommand{\I}{\mathbf{1}}
\newcommand{\Tr}{\operatorname{Tr}}
\newcommand{\rank}{\operatorname{rank}}
\newcommand{\Ran}{\operatorname{Ran}}
\newcommand{\diag}{\operatorname{diag}}
\newcommand{\sgn}{\operatorname{sgn}}
\newcommand{\ket}[1]{\lvert #1\rangle}
\newcommand{\bra}[1]{\langle #1\rvert}
\newcommand{\cH}{\mathcal{H}}
\newcommand{\cF}{\mathcal{F}}
\newcommand{\eps}{\varepsilon}

\theoremstyle{plain}
\newtheorem{proposition}{Proposition}

\theoremstyle{definition}
\newtheorem{remark}{Remark}

\begin{document}

\title{\texorpdfstring{\(U(1)\)}{U(1)}-Symmetric Fermionic Gaussian Transfer-Matrix Formalism}
\author{Asadullah Bhuiyan}
\date{\today}

\begin{abstract}
We organize the transfer-matrix formalism for number-conserving fermionic
Gaussian operators in the notation of the handwritten derivation.  A physical
Gaussian state is labeled by the complex covariance
\(G=2C-\I\), while a Gaussian operator is represented by a pure Choi covariance
\(\Sigma\) on the doubled one-particle space
\(\cH_L\oplus\cH_R\).  For a full-rank even Gaussian operator the doubled Choi
correlation matrix gives an ordinary particle transfer matrix
\(T_p=C_{LL}^{-1}C_{LR}\), with the corresponding hole transfer matrix
\(T_h=(T_p^\dagger)^{-1}\).  We then take the projective rank-one limits
needed for occupation projectors and single-fermion gain/loss branches.  The
central correction is that branch-resolved single-fermion insertions remain
pure Gaussian states, but their postselected Choi covariances are not
determined by the singular effective transfer matrix alone.  A right-leg
creation or annihilation branch gives \(T_{\rm eff}=T_p(1-P_\chi)\); the full
singular datum also records the null-sector occupation, filled for
\(\hat\chi^\dagger\) and empty for \(\hat\chi\).
\end{abstract}

\maketitle
\setcounter{tocdepth}{2}
\tableofcontents

\section{Hamiltonian-to-Covariance Map}

Let \(c_i,c_i^\dagger\), \(i=1,\ldots,N\), denote number-conserving fermion
modes on the one-particle Hilbert space \(\cH\simeq\mathbb C^N\).  For a
single-particle Hamiltonian \(h=h^\dagger\), write
\begin{equation}
  h=u\,\diag(\epsilon_1,\ldots,\epsilon_N)\,u^\dagger ,
  \qquad
  a_n^\dagger=\sum_i u_{in}c_i^\dagger .
  \label{eq:h-diagonalization}
\end{equation}
The thermal Gaussian density matrix
\begin{equation}
  \rho_h=\frac{1}{Z}\exp\!\left(-\sum_{ij}c_i^\dagger h_{ij}c_j\right)
  =
  \frac{1}{Z}\exp\!\left(-\sum_n\epsilon_n a_n^\dagger a_n\right)
  \label{eq:rho-h}
\end{equation}
has one-body correlation matrix
\begin{equation}
  C_{ij}=\Tr(\rho_h\,c_i^\dagger c_j),
  \qquad
  C=u\,f(h_{\rm diag})\,u^\dagger,
  \qquad
  f(\epsilon)=\frac{1}{1+e^\epsilon}.
  \label{eq:thermal-C}
\end{equation}
The complex covariance used below is
\begin{equation}
  G=2C-\I .
  \label{eq:G-convention}
\end{equation}
In the zero-temperature or fixed-charge Slater limit \(C\) is an orthogonal
projector and
\begin{equation}
  G^2=\I .
  \label{eq:pure-G}
\end{equation}
Equivalently, if the occupied orbitals form the columns of an isometry
\(U_{\rm occ}\), then
\begin{equation}
  \ket{G}
  =
  \prod_{\alpha=1}^{q}
  \left(\sum_i (U_{\rm occ})_{i\alpha}c_i^\dagger\right)\ket{0},
  \qquad
  C=U_{\rm occ}U_{\rm occ}^\dagger,
  \qquad
  G=2U_{\rm occ}U_{\rm occ}^\dagger-\I .
  \label{eq:slater-G}
\end{equation}
This is the state-covariance convention of the first part of the handwritten
notes.

\section{Particle and Hole Excitation Relations}

For a normalized one-particle mode
\begin{equation}
  \ket{\chi}=\sum_i\chi_i\ket{i},
  \qquad
  \hat\chi^\dagger=\sum_i\chi_i c_i^\dagger,
  \qquad
  \hat\chi=\sum_i\chi_i^*c_i,
  \qquad
  P_\chi=\ket{\chi}\bra{\chi},
  \label{eq:chi-mode}
\end{equation}
the only component that changes a fixed-charge Slater state is the component
orthogonal to the already occupied or empty subspace.  Thus, if the branch
amplitude is nonzero, particle addition and particle removal update the
physical state covariance by
\begin{align}
  C_{\rm add}
  &=
  C+
  \frac{(1-C)P_\chi(1-C)}
       {\Tr[(1-C)P_\chi]},
  \label{eq:state-addition}\\
  C_{\rm rem}
  &=
  C-
  \frac{CP_\chi C}
       {\Tr[CP_\chi]} .
  \label{eq:state-removal}
\end{align}
These equations are the fixed-charge state version of the Choi rank-one
updates derived later.  They also make clear which denominator is a Born
probability: a creation branch requires an empty component of \(\chi\), while
a removal branch requires an occupied component.

Let \(V\) be an invertible, number-conserving Gaussian operator.  Up to a scalar
it may be written as
\begin{equation}
  V=\exp\!\left(\sum_{ij}c_i^\dagger A_{ij}c_j\right).
  \label{eq:gaussian-operator}
\end{equation}
The particle and hole transfer matrices are defined by the conjugation rules
\begin{align}
  V\,c^\dagger(\phi)\,V^{-1}
  &=
  c^\dagger(T_p\phi),
  \label{eq:Tp-def}\\
  V\,c(\phi)\,V^{-1}
  &=
  c(T_h\phi),
  \label{eq:Th-def}
\end{align}
where \(c^\dagger(\phi)=\sum_i\phi_i c_i^\dagger\) and
\(c(\phi)=\sum_i\phi_i^*c_i\).  With the convention
\eqref{eq:gaussian-operator},
\begin{equation}
  T_p=e^A,
  \qquad
  T_h=(T_p^\dagger)^{-1}.
  \label{eq:Tp-Th-relation}
\end{equation}
If \(T_p=e^H u\) is the polar decomposition with \(H=H^\dagger\) and \(u\)
unitary, then
\begin{equation}
  T_h=e^{-H}u .
  \label{eq:polar-Th}
\end{equation}
This is the precise form of the particle/hole relation used in the transfer
matrix part of the handwritten derivation.

\section{Full-Rank \texorpdfstring{\(U(1)\)}{U(1)} Gaussian Operators}

The full-rank even chart consists of operators generated by fermion bilinears,
as in Eq.~\eqref{eq:gaussian-operator}, with \(T_p\) invertible.  Acting on a
Slater determinant, \(V\) transports the occupied one-particle frame by
\(T_p\), followed by the Gram normalization required to return to a normalized
Slater state.  Thus the physical state transform is nonlinear in \(C\), while
the excitation transform is linear in the one-particle transfer matrix.

The two transfer matrices \(T_p\) and \(T_h\) are not independent.  Equation
\eqref{eq:Tp-Th-relation} states that an invertible particle transfer matrix
already fixes the hole transfer matrix.  Conversely, if one begins from
\(T_h\), then
\begin{equation}
  T_p=(T_h^\dagger)^{-1}.
  \label{eq:Th-to-Tp}
\end{equation}
The singular cases below should be understood as limits of this full-rank
chart.  That distinction is important: a finite full-rank \(T_p\) is a complete
coordinate for an even Gaussian Choi covariance, but a singular effective
matrix is not a complete coordinate for an odd postselected Choi covariance.

\section{Choi Vectorization and Doubled Covariance Blocks}

Introduce two copies of the one-particle Hilbert space,
\(\cH_L\oplus\cH_R\), with fermion operators \(c_{Li}\) and \(c_{Ri}\).  The
reference state is the number-conserving maximally entangled state
\begin{equation}
  \ket{\Omega}_{LR}
  =
  2^{-N/2}\prod_{i=1}^{N}
  \left(c_{Li}^\dagger+c_{Ri}^\dagger\right)\ket{0}_{LR}.
  \label{eq:reference-state}
\end{equation}
For one mode this is
\begin{equation}
  \ket{\Omega}
  =
  \frac{\ket{1_L0_R}+\ket{0_L1_R}}{\sqrt2}.
  \label{eq:one-mode-reference}
\end{equation}
We vectorize an operator by acting on the right leg,
\begin{equation}
  \ket{V}\!\rangle
  =
  V_R\ket{\Omega}_{LR}.
  \label{eq:choi-vector}
\end{equation}
The resulting Gaussian Choi state is labeled by its doubled correlation matrix
and covariance
\begin{equation}
  C_{\alpha\beta,ij}
  =
  \frac{\langle\!\bra{V}c_{\alpha i}^\dagger c_{\beta j}\ket{V}\!\rangle}
       {\langle\!\bra{V}V\rangle\!\rangle},
  \qquad
  \Sigma=2C-\I_{L\oplus R},
  \qquad
  \alpha,\beta\in\{L,R\}.
  \label{eq:choi-C-Sigma}
\end{equation}
We use the block notation
\begin{equation}
  C=
  \begin{pmatrix}
    C_{LL} & C_{LR}\\
    C_{RL} & C_{RR}
  \end{pmatrix},
  \qquad
  \Sigma=
  \begin{pmatrix}
    \Sigma_{LL} & \Sigma_{LR}\\
    \Sigma_{RL} & \Sigma_{RR}
  \end{pmatrix}.
  \label{eq:block-notation}
\end{equation}
For a nonzero branch-resolved Gaussian Choi vector,
\begin{equation}
  C^2=C,
  \qquad
  \Sigma^2=\I_{L\oplus R}.
  \label{eq:choi-purity}
\end{equation}

\section{Full-Rank Transfer-Matrix/Choi Dictionary}

For an even full-rank Gaussian operator with particle transfer matrix
\(T\equiv T_p\), define
\begin{equation}
  L=(\I+TT^\dagger)^{-1},
  \qquad
  R=(\I+T^\dagger T)^{-1}.
  \label{eq:L-R-def}
\end{equation}
The Choi correlation blocks are
\begin{align}
  C_{LL}&=L,
  &
  C_{LR}&=LT=TR,
  \label{eq:choi-blocks-LL-LR}\\
  C_{RL}&=T^\dagger L=RT^\dagger,
  &
  C_{RR}&=T^\dagger LT=\I-R .
  \label{eq:choi-blocks-RL-RR}
\end{align}
Equivalently, in covariance variables,
\begin{align}
  \Sigma_{LL}&=2L-\I,
  &
  \Sigma_{LR}&=2LT,
  \label{eq:sigma-blocks-LL-LR}\\
  \Sigma_{RL}&=2T^\dagger L,
  &
  \Sigma_{RR}&=\I-2R .
  \label{eq:sigma-blocks-RL-RR}
\end{align}
The inverse map in the full-rank chart is
\begin{equation}
  T=C_{LL}^{-1}C_{LR}
  =
  (\I+\Sigma_{LL})^{-1}\Sigma_{LR}.
  \label{eq:T-extraction}
\end{equation}
The second equality follows from \(\I+\Sigma_{LL}=2C_{LL}\) and
\(\Sigma_{LR}=2C_{LR}\).  Equation \eqref{eq:T-extraction} is the formula that
must be used only on the support where the even full-rank chart is valid.

\begin{proposition}
The blocks \eqref{eq:choi-blocks-LL-LR} and
\eqref{eq:choi-blocks-RL-RR} define a pure Choi state:
\(C^2=C\), or equivalently \(\Sigma^2=\I\).
\end{proposition}

\begin{proof}
The identities \(LT=TR\), \(LTT^\dagger=\I-L\), and
\(RT^\dagger T=\I-R\) imply
\begin{align}
  C_{LL}^2+C_{LR}C_{RL}
  &=
  L^2+LTT^\dagger L
  =
  L,\\
  C_{LL}C_{LR}+C_{LR}C_{RR}
  &=
  L^2T+LT(\I-R)
  =
  LT,
\end{align}
with the remaining two block equations following by Hermitian conjugation and
the same \(R\)-identity.
\end{proof}

\section{Rank-One Projectors and Projective Limits}

Let \(P=P_\chi=\ket{\chi}\bra{\chi}\) and \(Q=\I-P\).  The occupation
projectors \(N_\chi=\hat\chi^\dagger\hat\chi\) and \(1-N_\chi\) arise as
projective limits of full-rank even Gaussian operators:
\begin{align}
  N_\chi
  &=
  \lim_{\lambda\to\infty}
  e^{-\lambda}e^{\lambda N_\chi},
  &
  T_\lambda^{(N)}
  &=
  Q+e^\lambda P,
  \label{eq:N-projective}\\
  1-N_\chi
  &=
  \lim_{\lambda\to\infty}
  e^{-\lambda N_\chi},
  &
  T_\lambda^{(1-N)}
  &=
  Q+e^{-\lambda}P .
  \label{eq:one-minus-N-projective}
\end{align}
At the transfer-matrix level the first limit diverges on \(\Ran P\), while the
second becomes rank deficient.  At the Choi-covariance level both limits are
finite.  On the \(P\)-sector,
\begin{equation}
  C(T=t)
  =
  \frac{1}{1+|t|^2}
  \begin{pmatrix}
    1 & t\\
    t^* & |t|^2
  \end{pmatrix}_{LR}.
  \label{eq:one-mode-even-C}
\end{equation}
Thus \(t\to\infty\) gives
\(\ket{0_L1_R}\), whereas \(t\to0\) gives \(\ket{1_L0_R}\).  These are the
two even null completions of the full-rank chart.  They are not the odd
single-fermion gain/loss completions discussed below.

\section{Single-Fermion Gain and Loss Branches}

Now let \(V\) be an even full-rank Gaussian operator and consider the
branch-resolved odd Kraus amplitudes
\begin{equation}
  K_+=\hat\chi^\dagger V,
  \qquad
  K_-=\hat\chi V .
  \label{eq:odd-Kraus}
\end{equation}
Because \(K_\pm\) are odd, they do not define ordinary invertible even
conjugation maps.  The correct object is instead the Choi vector
\begin{equation}
  \ket{K_+}\!\rangle=\hat\chi_R^\dagger\ket{V}\!\rangle,
  \qquad
  \ket{K_-}\!\rangle=\hat\chi_R\ket{V}\!\rangle.
  \label{eq:odd-choi-vectors}
\end{equation}
If the corresponding branch amplitude is nonzero, these are still pure
Gaussian Choi states.  Therefore the updated covariance still satisfies
\begin{equation}
  \Sigma_\pm^2=\I .
  \label{eq:odd-purity}
\end{equation}
The failure is not purity.  The failure is the assumption that an odd
postselected Choi state can be reconstructed from the ordinary full-rank even
transfer coordinate alone.

Set
\begin{equation}
  u_R=
  \begin{pmatrix}
    0\\
    \ket{\chi}
  \end{pmatrix}_{L\oplus R}.
  \label{eq:uR}
\end{equation}
Creation fills the empty projection of \(u_R\), while annihilation removes its
occupied projection:
\begin{align}
  C_+
  &=
  C+
  \frac{(\I-C)u_Ru_R^\dagger(\I-C)}
       {u_R^\dagger(\I-C)u_R},
  \label{eq:C-plus-update}\\
  C_-
  &=
  C-
  \frac{Cu_Ru_R^\dagger C}
       {u_R^\dagger C u_R}.
  \label{eq:C-minus-update}
\end{align}
The denominators
\begin{equation}
  d_+
  =
  \bra{\chi}(\I-C_{RR})\ket{\chi}
  =
  \Tr[(\I-C_{RR})P],
  \qquad
  d_-
  =
  \bra{\chi}C_{RR}\ket{\chi}
  =
  \Tr[C_{RR}P]
  \label{eq:d-plus-minus}
\end{equation}
are the branch probabilities in the normalized Choi state.  The formulas are
defined for \(d_\pm\neq0\).

Blockwise, the creation branch is
\begin{align}
  C_{LL}^{(+)}
  &=
  C_{LL}
  +
  \frac{C_{LR}PC_{RL}}{d_+},
  \label{eq:C-plus-LL}\\
  C_{LR}^{(+)}
  &=
  C_{LR}
  -
  \frac{C_{LR}P(\I-C_{RR})}{d_+},
  \label{eq:C-plus-LR}\\
  C_{RR}^{(+)}
  &=
  C_{RR}
  +
  \frac{(\I-C_{RR})P(\I-C_{RR})}{d_+},
  \label{eq:C-plus-RR}
\end{align}
with \(C_{RL}^{(+)}=(C_{LR}^{(+)})^\dagger\).  The loss branch is
\begin{align}
  C_{LL}^{(-)}
  &=
  C_{LL}
  -
  \frac{C_{LR}PC_{RL}}{d_-},
  \label{eq:C-minus-LL}\\
  C_{LR}^{(-)}
  &=
  C_{LR}
  -
  \frac{C_{LR}PC_{RR}}{d_-},
  \label{eq:C-minus-LR}\\
  C_{RR}^{(-)}
  &=
  C_{RR}
  -
  \frac{C_{RR}PC_{RR}}{d_-},
  \label{eq:C-minus-RR}
\end{align}
with \(C_{RL}^{(-)}=(C_{LR}^{(-)})^\dagger\).

\section{Effective Singular Transfer Matrix}

Suppose the pre-insertion Choi covariance is in the full-rank even chart
\eqref{eq:choi-blocks-LL-LR}--\eqref{eq:choi-blocks-RL-RR}, with
\(T=T_p\).  After the right-leg insertion the ordinary inverse in
\eqref{eq:T-extraction} must be replaced by the support inverse, or
equivalently by a full-rank regularization followed by the singular limit.
Substituting Eqs.~\eqref{eq:choi-blocks-LL-LR}--\eqref{eq:choi-blocks-RL-RR}
into the rank-one updates gives
\begin{equation}
  T_{\rm eff}^{(+)}
  =
  (C_{LL}^{(+)})^{-1}_{\rm supp}C_{LR}^{(+)}
  =
  TQ,
  \qquad
  Q=\I-P ,
  \label{eq:Teff-plus}
\end{equation}
and similarly
\begin{equation}
  T_{\rm eff}^{(-)}
  =
  (C_{LL}^{(-)})^{-1}_{\rm supp}C_{LR}^{(-)}
  =
  TQ .
  \label{eq:Teff-minus}
\end{equation}
Thus a right-leg single-fermion insertion removes the \(\chi\)-column
direction of the particle transfer matrix:
\begin{equation}
  T_p\longmapsto T_p(\I-P_\chi).
  \label{eq:right-leg-column-removal}
\end{equation}
The analogous left-leg insertion removes the corresponding row direction:
\begin{equation}
  T_p\longmapsto(\I-P_\chi)T_p .
  \label{eq:left-leg-row-removal}
\end{equation}

\begin{remark}[First corrected pitfall: regulating the final odd covariance]
One may regulate the parent even operator or use the exact rank-one update
\eqref{eq:C-plus-update} or \eqref{eq:C-minus-update}.  What is not legitimate
is to interpolate the final odd covariance by hand and then continue to impose
pure-state transfer identities.  For example,
\begin{equation}
  \Sigma_\eps=
  \begin{pmatrix}
    P+\eps Q & Q+\eps P\\
    Q+\eps P & P+\eps Q
  \end{pmatrix}
  \label{eq:bad-regulator}
\end{equation}
has, on the \(P\)-sector,
\begin{equation}
  \Sigma_\eps\big|_P=
  \begin{pmatrix}
    1 & \eps\\
    \eps & 1
  \end{pmatrix},
  \qquad
  \Sigma_\eps^2\big|_P=
  \begin{pmatrix}
    1+\eps^2 & 2\eps\\
    2\eps & 1+\eps^2
  \end{pmatrix}
  \neq\I .
  \label{eq:bad-regulator-square}
\end{equation}
Hence such a finite-\(\eps\) object is not a pure Gaussian Choi covariance.
\end{remark}

\section{Null-Sector Occupation Data}

The effective matrix \(T_{\rm eff}=TQ\) captures the homogeneous propagation
sector, but it does not determine the full postselected Choi covariance.  It
only says that the right-leg \(P\)-column no longer participates in the
\(L\)-to-\(R\) bridge.  It does not say whether that disconnected mode is empty
or filled.

The distinction is already visible for one mode with \(T=1\).  The reference
state is Eq.~\eqref{eq:one-mode-reference}.  Right-leg creation gives
\begin{equation}
  \hat\chi_R^\dagger\ket{\Omega}
  \propto
  \ket{1_L1_R},
  \qquad
  C_+
  =
  \begin{pmatrix}
    1 & 0\\
    0 & 1
  \end{pmatrix}_{LR},
  \label{eq:one-mode-creation}
\end{equation}
whereas right-leg loss gives
\begin{equation}
  \hat\chi_R\ket{\Omega}
  \propto
  \ket{0_L0_R},
  \qquad
  C_-
  =
  \begin{pmatrix}
    0 & 0\\
    0 & 0
  \end{pmatrix}_{LR}.
  \label{eq:one-mode-loss}
\end{equation}
Both branches have \(T_{\rm eff}=0\).  They are nevertheless different pure
Choi states.

Equivalently, reconstructing from \(T_{\rm eff}=0\) using only the even formula
\eqref{eq:one-mode-even-C} gives one of the even null completions:
\begin{equation}
  t\to0:
  \quad
  C(t)\to
  \begin{pmatrix}
    1 & 0\\
    0 & 0
  \end{pmatrix}_{LR}
  \quad
  (\ket{1_L0_R}),
  \label{eq:even-null-tzero}
\end{equation}
or
\begin{equation}
  t\to\infty:
  \quad
  C(t)\to
  \begin{pmatrix}
    0 & 0\\
    0 & 1
  \end{pmatrix}_{LR}
  \quad
  (\ket{0_L1_R}).
  \label{eq:even-null-tinfty}
\end{equation}
Neither is the odd creation completion \(\ket{1_L1_R}\), and neither is the
odd loss completion \(\ket{0_L0_R}\).

\begin{remark}[Second corrected pitfall: regularizing \(T_{\rm eff}\) alone]
A singular effective transfer matrix can be regularized as an even matrix, but
that regularization chooses an even null completion.  It cannot recover which
odd branch occurred.  The missing information is precisely the null-sector
occupation.
\end{remark}

The complete singular data for a right-leg odd branch are therefore
\begin{equation}
  \left(T_{\rm eff},P_{\rm null},n_{\rm null}\right),
  \qquad
  P_{\rm null}=P_\chi,
  \label{eq:singular-data}
\end{equation}
where the occupation datum may be read from the right-right Choi block:
\begin{equation}
  P_\chi C_{RR}^{(+)}P_\chi=P_\chi,
  \qquad
  P_\chi C_{RR}^{(-)}P_\chi=0.
  \label{eq:null-occupation}
\end{equation}
Thus
\begin{equation}
  n_{\rm null}=
  \begin{cases}
    1, & \text{right-leg creation } \hat\chi_R^\dagger,\\
    0, & \text{right-leg loss } \hat\chi_R .
  \end{cases}
  \label{eq:n-null}
\end{equation}

\section{Final Picture}

The full-rank even formalism is internally consistent and complete:
\begin{equation}
  T_p
  \quad\Longleftrightarrow\quad
  C(T_p)
  \quad\Longleftrightarrow\quad
  \Sigma(T_p),
  \qquad
  T_h=(T_p^\dagger)^{-1}.
  \label{eq:even-equivalence}
\end{equation}
Rank-one even projectors lie on the boundary of this chart: the transfer
matrix may diverge or lose rank, but the Choi covariance has a finite
projective limit.  Odd single-fermion branches go one step further.  They are
valid pure Gaussian Choi states, and their covariances must be computed by the
rank-one creation/loss updates \eqref{eq:C-plus-update} and
\eqref{eq:C-minus-update}.  The extracted singular transfer matrix is
\begin{equation}
  T_{\rm eff}=T_p(\I-P_\chi)
  \label{eq:final-Teff}
\end{equation}
for a right-leg insertion, but this matrix is only the propagation part of the
answer.  The full postselected Choi covariance also requires the null-sector
occupation in Eq.~\eqref{eq:null-occupation}.  This resolves the apparent
contradiction in the handwritten derivation: purity is preserved, the Choi
representation remains valid, and the effective transfer matrix remains useful,
but the singular odd branch is not coordinatized by \(T_{\rm eff}\) alone.

\end{document}
